\clearpage
\chapter*{Fe de erratas de la exposición n.º 221}
\addcontentsline{toc}{chapter}{Fe de erratas de la exposición n.º 221}
\setcounter{footnote}{3}
\src{302}

\noindent\textbf{6. TDTE IV: Los esquemas de Hilbert.}

\begin{center}
[Seminario Bourbaki, t. 13, 1960/61, n.º 221, 28 pp.]
\end{center}

\noindent\textbf{Página 6, Teorema 2.1, 3.ª línea.} --- En lugar de
«es necesario», léase «es necesario y suficiente».

\noindent\textbf{Página 15, línea 13.} --- En lugar de «$P\neq Q$», léase
«$P<Q$».

\noindent\textbf{Línea 15.} --- En lugar de «Si, por el contrario, $P=Q$»,
léase «En el caso contrario, se tendrá $P(n)\geq Q(n)$ para $n$ grande».

\noindent\textbf{Línea 18, fórmula (*).} --- En lugar de «$<$» y «$>$»,
léase «$\leq$» y «$\geq$».

\noindent\textbf{5.ª línea desde abajo.} --- En lugar de «$r$», léase «$Z$».

\noindent\textbf{Página 18, Observación 3.9.} --- Señalemos que Hartshorne ha
estudiado las componentes conexas de los esquemas de Hilbert sobre un cuerpo
algebraicamente cerrado. Este autor demuestra que los
$\operatorname{Hilb}^{P}_{r}$ son conexos y determina los pares $(r,P)$ para
los cuales $\operatorname{Hilb}^{P}_{r}\neq\varnothing$\footnote{HARTSHORNE
(R.). --- Connectedness of the Hilbert scheme (Tesis, Harvard, 1962).}.

\noindent\textbf{Página 23, 5.ª línea desde abajo.} --- En lugar de «una
componente», léase «de una componente».

\noindent\textbf{Páginas 23 y 24, Observación 5.5.} --- A propósito del ejemplo
de Zappa, señalemos que Mumford acaba de construir una componente irreducible
del esquema de Hilbert de $\mathbf P^3_{\mathbf C}$ (cuyos puntos generales
representan curvas no singulares de grado 14 y género 24) que es no reducida
en su punto genérico. Haciendo explotar una de las curvas obtenidas, obtiene
también un esquema proyectivo regular de dimensión 3 sobre $\mathbf C$, cuyo
esquema formal de módulos es no reducido en su punto genérico o, lo que
equivale a lo mismo, cuya variedad local de módulos, en el sentido de la
geometría analítica sobre $\mathbf C$ (cf. Seminario Cartan, t. 13, 1960/61),
es no reducida en todos sus puntos.
